%% toc: P3 31 31.9
\section{Proofs for self-extending models and indiscernibles of arbitrary length}
\label{sec:annex-p03-w03-p03-030-self-extending}

Throughout, \(L\) is a first-order vocabulary and formulas are formulas of
\(\mathscr L(Q)\) unless they are called first-order.  For a structure
\((\mathfrak A,q)\), a formula \(\phi(x,\bar y)\) and parameters \(\bar b\),
write \(\phi(A,\bar b)\) for \(\{a\in A:(\mathfrak A,q)\entails\phi(a,\bar
b)\}\).  So \((\mathfrak A,q)\entails Qx\,\phi(x,\bar b)\) means
\(\phi(A,\bar b)\in q\).

\begin{annexitem}{annex:p03-w03-p03-030-ultrafilter-facts}{thm:p03-w03-p03-030-ultrafilter-facts}{The definable sets in \(q\) form a nonprincipal ultrafilter}
\begin{proof}
Satisfaction of \(Qx\,\phi\) depends only on the set \(\phi(A,\bar b)\), so
each scheme instance is a statement about definable sets.

\emph{Complements.}  The first scheme says that \(X\in q\) iff
\(A\setminus X\notin q\), for every definable \(X\), since
\((\neg\phi)(A,\bar b)=A\setminus\phi(A,\bar b)\).  So exactly one of \(X\)
and \(A\setminus X\) is in \(q\).

\emph{Supersets and intersections.}  If \(X=\phi(A,\bar b)\subseteq
Y=\psi(A,\bar c)\), then \(\forall x\,(\phi(x,\bar b)\rightarrow\psi(x,\bar
c))\) holds, and the third scheme gives \(X\in q\Rightarrow Y\in q\).  The
fourth scheme gives \(X,Y\in q\Rightarrow X\cap Y\in q\).

\emph{\(A\in q\) and \(\varnothing\notin q\).}  If \(\varnothing\in q\), pick
\(b\in A\).  Then \(\varnothing\subseteq\{b\}\), and \(\{b\}\) is definable
by \(x=b\), so \(\{b\}\in q\), contradicting the scheme
\(\forall y\,\neg Qx\,(x=y)\).  Thus \(\varnothing\notin q\), and by
complements \(A\in q\).  This proves (1).

(2) The second scheme says that no singleton is in \(q\).  If \(A\) were
finite, \(A=\{b_1,\ldots,b_k\}\), then \(A\in q\) and the complement rule
applied to \(\{b_1\}\), then to \(\{b_2\}\) inside \(A\setminus\{b_1\}\),
and so on, would give: since \(\{b_1\}\notin q\), \(A\setminus\{b_1\}\in q\);
intersecting with \(A\setminus\{b_2\}\in q\) gives
\(A\setminus\{b_1,b_2\}\in q\); after \(k\) steps \(\varnothing\in q\), a
contradiction.  So \(A\) is infinite.

(3) If \(\phi(A,\bar b)\in q\) then \(\phi(A,\bar b)\ne\varnothing\) by (1),
so \(\exists x\,\phi(x,\bar b)\) holds.  If \(x\) is not free in \(\chi\),
then \(\chi(A,\bar b)\) is \(A\) when \(\chi(\bar b)\) holds and
\(\varnothing\) otherwise, and (1) gives \(Qx\,\chi\leftrightarrow\chi\).
\end{proof}
\end{annexitem}

\begin{annexitem}{annex:p03-w03-p03-030-compactness}{thm:p03-w03-p03-030-compactness}{Compactness for structures for \(\mathscr L(Q)\)}
\begin{proof}
\emph{The translation.}  For every \(\mathscr L(Q)\)-formula of the form
\(Qx\,\phi\), list its free variables in a fixed order \(\bar y\) and add to
\(L\) a new relation symbol \(P_{Qx\phi}\) of arity the length of \(\bar y\).
Call the enlarged vocabulary \(L^\sharp\).  Translate each
\(\mathscr L(Q)\)-formula \(\chi\) into a first-order \(L^\sharp\)-formula
\(\chi^\sharp\) with the same free variables: atomic formulas are unchanged,
the translation commutes with \(\neg\), \(\wedge\) and \(\exists y\), and
\((Qx\,\phi)^\sharp\defeq P_{Qx\phi}(\bar y)\).  Let \(\mathrm{Ext}\) be the
set of universal closures of all formulas
\[
 \forall x\,\bigl(\phi^\sharp(x,\bar y)\leftrightarrow\psi^\sharp(x,\bar z)\bigr)
 \rightarrow\bigl(P_{Qx\phi}(\bar y)\leftrightarrow P_{Qx'\psi'}(\bar z)\bigr),
\]
where \(Qx\,\phi(x,\bar y)\) and \(Qx'\,\psi'(x',\bar z)\) are any two such
formulas whose free-variable lists \(\bar y\) and \(\bar z\) have been renamed
apart, \(x\) is a variable occurring in neither list, and \(\psi(x,\bar z)\)
is \(\psi'\) with \(x\) put for its bound variable \(x'\).

\emph{From \(\mathscr L(Q)\)-structures to \(L^\sharp\)-structures.}  Given
\((\mathfrak A,q)\), expand \(\mathfrak A\) to \(\mathfrak A^\sharp\) by
letting \(P_{Qx\phi}(\bar b)\) hold iff \((\mathfrak A,q)\entails
Qx\,\phi(x,\bar b)\).  By induction on \(\chi\),
\((\mathfrak A,q)\entails\chi(\bar a)\) iff \(\mathfrak
A^\sharp\entails\chi^\sharp(\bar a)\): the atomic and Boolean steps and the
step for \(\exists\) are the same in both structures, and the step for \(Q\)
holds by the definition of \(P_{Qx\phi}\).  \(\mathfrak A^\sharp\) satisfies
\(\mathrm{Ext}\).  If the antecedent holds at \(\bar b,\bar c\), then by the
equivalence just proved \(\{a:(\mathfrak A,q)\entails\phi(a,\bar b)\}\) and
\(\{a:(\mathfrak A,q)\entails\psi(a,\bar c)\}\) are the same set, and the
truth of \(Qx\,\phi(x,\bar b)\) and of \(Qx'\,\psi'(x',\bar c)\) depends only
on whether that set is in \(q\).

\emph{From \(L^\sharp\)-structures to \(\mathscr L(Q)\)-structures.}  Let
\(\mathfrak C\) be an \(L^\sharp\)-structure satisfying \(\mathrm{Ext}\), with
\(L\)-reduct \(\mathfrak A\).  Put
\[
 q\defeq\{\phi^\sharp(\mathfrak C,\bar b):\mathfrak C\entails
 P_{Qx\phi}(\bar b)\},
 \qquad
 \phi^\sharp(\mathfrak C,\bar b)\defeq\{a:\mathfrak C\entails\phi^\sharp(a,\bar b)\}.
\]
We show by induction on \(\chi\) that \((\mathfrak A,q)\entails\chi(\bar a)\)
iff \(\mathfrak C\entails\chi^\sharp(\bar a)\).  Only the step for \(Q\)
needs comment.  By the induction hypothesis the set defined by
\(\phi(x,\bar b)\) in \((\mathfrak A,q)\) is \(\phi^\sharp(\mathfrak
C,\bar b)\).  If \(\mathfrak C\entails P_{Qx\phi}(\bar b)\), that set is in
\(q\) by definition.  Conversely, if it is in \(q\), there are
\(Qx'\,\psi'\) and \(\bar c\) with \(\mathfrak C\entails
P_{Qx'\psi'}(\bar c)\) and \(\psi^\sharp(\mathfrak C,\bar c)=\phi^\sharp(\mathfrak
C,\bar b)\).  The instance of \(\mathrm{Ext}\) for this pair gives
\(\mathfrak C\entails P_{Qx\phi}(\bar b)\).

\emph{Compactness.}  Let \(\Gamma^\sharp=\{\gamma^\sharp:\gamma\in\Gamma\}\).
A finite subset of \(\Gamma^\sharp\cup\mathrm{Ext}\) contains the
translations of finitely many members of \(\Gamma\).  Those have a model
\((\mathfrak A,q)\), and \(\mathfrak A^\sharp\) satisfies their translations
and all of \(\mathrm{Ext}\).  By the first-order compactness theorem
(\cref{thm:fol-compactness}), which uses the axiom of choice when \(L^\sharp\)
is uncountable, \(\Gamma^\sharp\cup\mathrm{Ext}\) has a model \(\mathfrak C\).
The structure \((\mathfrak A,q)\) built from \(\mathfrak C\) is then a model
of \(\Gamma\).
\end{proof}
\end{annexitem}

\begin{annexitem}{annex:p03-w03-p03-030-canonical-extension}{thm:p03-w03-p03-030-canonical-extension}{The canonical extension and its \L o\'s equivalence}
\begin{proof}
For definable \(f_0,\ldots,f_{n-1}\) and a formula \(\psi(y_0,\ldots,y_{n-1})\)
put \(\psi[\bar f]\defeq\{a\in A:(\mathfrak A,q)\entails\psi(f_0(a),\ldots,
f_{n-1}(a))\}\).  This set is definable, by the formula in
\cref{def:p03-w03-p03-030-canonical-extension}.  If \(g\) is definable and
\(\tau(\bar y)\) is an \(L\)-term, the function
\(x\mapsto\tau^{\mathfrak A}(f_0(x),\ldots,f_{n-1}(x))\) is definable,
because its graph is defined by \(\exists\bar y(\bigwedge_i\gamma_i(x,y_i)\wedge
z=\tau(\bar y))\).  All uses of ``in \(q\)'' below rely on
\cref{thm:p03-w03-p03-030-ultrafilter-facts}.

(1) \emph{Well-definedness.}  Equality modulo \(q\) is reflexive because
\(A\in q\).  It is symmetric by definition, and transitive because
\(\{f=g\}\cap\{g=h\}\subseteq\{f=h\}\).  Suppose \(f_i\) and \(f'_i\) are
equal modulo \(q\) for each \(i<n\), and let \(Y=\bigcap_{i<n}\{a:f_i(a)=
f'_i(a)\}\in q\).  For any formula \(\psi\),
\(Y\cap\psi[\bar f]\subseteq\psi[\bar f']\) and
\(Y\cap\psi[\bar f']\subseteq\psi[\bar f]\).  Hence
\(\psi[\bar f]\in q\iff\psi[\bar f']\in q\).  With \(\psi=R(\bar y)\) this
shows that \(R^{\mathfrak B}\) is well defined.  With \(\psi\) the formula
\(F(\bar y)=z\), and the extra function \(x\mapsto F^{\mathfrak A}(\bar
f(x))\) as the last argument, it shows \([F(\bar f)]=[F(\bar f')]\).  In the
definition of a set in \(r\), the condition on \([g]\) is
\(\phi[\bar f,g]\in q\), which by the same argument does not depend on the
representative \(g\).  So each member of \(r\) is a well-defined set of
classes.

\emph{Terms.}  By induction on an \(L\)-term \(\tau(\bar y)\),
\(\tau^{\mathfrak B}([f_0],\ldots,[f_{n-1}])=[x\mapsto\tau^{\mathfrak A}(\bar
f(x))]\): variables and constants are immediate, and the step for a
function symbol is the definition of \(F^{\mathfrak B}\).

(2) \emph{The \L o\'s equivalence}, by induction on \(\phi\), for all tuples
of definable functions simultaneously.

\emph{Atomic.}  For \(R(\tau_1,\ldots,\tau_k)\), let
\(h_j(x)=\tau_j^{\mathfrak A}(\bar f(x))\).  By the term clause and the
definition of \(R^{\mathfrak B}\), \((\mathfrak B,r)\entails
R(\bar\tau)([\bar f])\) iff \(\{a:R^{\mathfrak A}(h_1(a),\ldots,h_k(a))\}\in
q\), which is \(R(\bar\tau)[\bar f]\in q\).  For \(\tau_1=\tau_2\), the classes
\([h_1]\) and \([h_2]\) are equal iff \(\{a:h_1(a)=h_2(a)\}\in q\), by the
definition of \(B\).

\emph{Negation.}  \((\mathfrak B,r)\entails\neg\phi([\bar f])\) iff
\(\phi[\bar f]\notin q\) (induction hypothesis), iff
\(A\setminus\phi[\bar f]=(\neg\phi)[\bar f]\in q\) (complements).

\emph{Conjunction.}  \((\phi\wedge\psi)[\bar f]=\phi[\bar f]\cap\psi[\bar
f]\), which is in \(q\) iff both sets are, by intersections and supersets.

\emph{Existential quantifier.}  Suppose \((\mathfrak B,r)\entails\exists
y\,\phi([\bar f],y)\).  Then for some definable \(g\),
\(\phi[\bar f,g]\in q\) by the induction hypothesis, and
\(\phi[\bar f,g]\subseteq(\exists y\,\phi)[\bar f]\), so the latter is in
\(q\).  Conversely, suppose \((\exists y\,\phi)[\bar f]\in q\).  By the Skolem
property there is a term \(t(\bar y)\) with
\(\forall\bar y\,(\exists y'\phi(\bar y,y')\rightarrow\phi(\bar
y,t(\bar y)))\in T\).  Let \(g(x)=t^{\mathfrak A}(\bar f(x))\), which is
definable.  Then \((\exists y\,\phi)[\bar f]\subseteq\phi[\bar f,g]\), so
\(\phi[\bar f,g]\in q\), and \((\mathfrak B,r)\entails\phi([\bar f],[g])\) by
the induction hypothesis.

\emph{The quantifier \(Q\).}  Let \(X=\{[g]\in B:(\mathfrak B,r)\entails
\phi([\bar f],[g])\}\).  By the induction hypothesis,
\(X=\{[g]:\phi[\bar f,g]\in q\}\), and \((\mathfrak B,r)\entails
Qy\,\phi([\bar f],y)\) iff \(X\in r\).  If \((Qy\,\phi)[\bar f]\in q\), that
is, \((\mathfrak A,q)\entails Qx\,Qy\,\phi(\bar f(x),y)\), then \(X\in r\) by
the definition of \(r\).  Conversely, let \(X\in r\).  Then
\(X=\{[g]:\psi[\bar h,g]\in q\}\) for some formula \(\psi\) and definable
\(\bar h\) with \((Qy\,\psi)[\bar h]\in q\).  Suppose, for a contradiction,
that \((Qy\,\phi)[\bar f]\notin q\).  By complements
\((\neg Qy\,\phi)[\bar f]\in q\).  In every model of \(T\),
\(\neg Qy\,\phi\) is equivalent to \(Qy\,\neg\phi\), by the first scheme.
So the set
\[
 Z=\{a:(\mathfrak A,q)\entails Qy\,\neg\phi(\bar f(a),y)\wedge
                              Qy\,\psi(\bar h(a),y)\}
\]
is in \(q\).  For each \(a\in Z\), the fourth scheme and part (3) of
\cref{thm:p03-w03-p03-030-ultrafilter-facts} give
\(\exists y\,(\neg\phi(\bar f(a),y)\wedge\psi(\bar h(a),y))\).  By the Skolem
property applied to the formula \(\neg\phi(\bar u,y)\wedge\psi(\bar v,y)\),
there is a term \(t\) such that \(g(x)=t^{\mathfrak A}(\bar f(x),\bar h(x))\)
satisfies \(Z\subseteq(\neg\phi)[\bar f,g]\cap\psi[\bar h,g]\).  Hence
\(\psi[\bar h,g]\in q\), so \([g]\in X\), and \((\neg\phi)[\bar f,g]\in q\), so
\(\phi[\bar f,g]\notin q\) and \([g]\notin X\).  This contradiction shows
\((Qy\,\phi)[\bar f]\in q\).

(3) \emph{Reducedness and uniqueness.}  A member of \(r\) has the form
\(\{[g]:\psi[\bar h,g]\in q\}\), which by (2) is
\(\{b\in B:(\mathfrak B,r)\entails\psi([\bar h],b)\}\), a definable set with
parameters \([\bar h]\).  So \((\mathfrak B,r)\) is reduced.  Let
\((\mathfrak B',r')\) be a reduced structure with domain \(B\) satisfying the
equivalence in (2).  Applying the equivalence to the formulas
\(R(\bar y)\), \(F(\bar y)=z\) and \(c=z\) shows that \(\mathfrak B'\) and
\(\mathfrak B\) interpret every symbol of \(L\) alike.  Let \(S\in r'\).
Since \((\mathfrak B',r')\) is reduced, \(S=\{b:(\mathfrak B',r')\entails
\psi([\bar h],b)\}\) for some formula \(\psi\) and some \([\bar h]\), and
\((\mathfrak B',r')\entails Qy\,\psi([\bar h],y)\).  By the equivalence,
\((Qy\,\psi)[\bar h]\in q\) and \(S=\{[g]:\psi[\bar h,g]\in q\}\), so
\(S\in r\).  Conversely, a member \(S=\{[g]:\psi[\bar h,g]\in q\}\) of \(r\)
with \((Qy\,\psi)[\bar h]\in q\) equals \(\{b:(\mathfrak B',r')\entails
\psi([\bar h],b)\}\), and \((\mathfrak B',r')\entails Qy\,\psi([\bar h],y)\)
by the equivalence, so \(S\in r'\).  Thus \(r'=r\).

(4) \emph{Elementary embedding.}  If \(a\ne a'\) then
\(\{x:c_a(x)=c_{a'}(x)\}=\varnothing\notin q\), so \(a\mapsto[c_a]\) is
injective.  For a formula \(\phi(\bar y)\) and \(\bar a\) from \(A\), (2) gives
\((\mathfrak B,r)\entails\phi([c_{\bar a}])\) iff \((\mathfrak A,q)\entails
Qx\,\phi(\bar a)\), with \(x\) not free.  By part (3) of
\cref{thm:p03-w03-p03-030-ultrafilter-facts}, this holds iff
\((\mathfrak A,q)\entails\phi(\bar a)\).  Applied to sentences, this shows
\((\mathfrak B,r)\entails T\).

(5) \emph{Cardinality.}  A definable function is determined by a formula and
a finite tuple of parameters from \(A\).  The set \(A\) is infinite by
\cref{thm:p03-w03-p03-030-ultrafilter-facts}, so there are at most
\((\abs L+\aleph_0)\cdot\abs A=\abs A+\abs L+\aleph_0\) definable functions.
Hence \(\abs B\le\abs A+\abs L+\aleph_0\), and \(\abs A\le\abs B\) by (4).
\end{proof}
\end{annexitem}

\begin{annexitem}{annex:p03-w03-p03-030-large-sets}{thm:p03-w03-p03-030-large-sets}{Large sets become larger}
\begin{proof}
Let \(\bar a\) be from \(A\), and suppose \((\mathfrak B,r)\entails
Qx\,\phi(x,\bar a)\).  By part (4) of
\cref{thm:p03-w03-p03-030-canonical-extension},
\((\mathfrak A,q)\entails Qx\,\phi(x,\bar a)\).  This says
\(\phi[\mathrm{id},c_{\bar a}]\in q\), where \(\mathrm{id}\) is the identity
function.  By part (2) of that theorem,
\((\mathfrak B,r)\entails\phi([\mathrm{id}],[c_{\bar a}])\), that is,
\(\phi(\iota,\bar a)\).  If \(\iota=[c_a]\) for some \(a\in A\), then
\(\{x:x=a\}\in q\), contradicting the scheme \(\forall y\,\neg Qx\,(x=y)\).
So \(\iota\notin A\).
\end{proof}
\end{annexitem}

\begin{annexitem}{annex:p03-w03-p03-030-hull}{thm:p03-w03-p03-030-hull}{Small models of a self-extending theory}
\begin{proof}
Let \((\mathfrak A,q)\) be a reduced model of \(T\), and fix \(a\in A\).  Let
\(H=\{s^{\mathfrak A}(a):s(v)\text{ an \(L\)-term in the one variable }v\}\).
Then \(H\) contains the constants and is closed under the functions of
\(\mathfrak A\), so it is the domain of a substructure \(\mathfrak H\).  It has
at most \(\abs L+\aleph_0\) elements.  If \(\bar h\) is from \(H\) and \(t\)
is an \(L\)-term, then \(t^{\mathfrak A}(\bar h)\in H\), since substituting
terms in \(v\) into \(t\) gives a term in \(v\).

Put \(q_H\defeq\{\phi(A,\bar h)\cap H:\bar h\text{ from }H,\
\phi(A,\bar h)\in q\}\).  We show, by induction on \(\chi\), that
\((\mathfrak H,q_H)\entails\chi(\bar h)\) iff \((\mathfrak A,q)\entails
\chi(\bar h)\) for all \(\bar h\) from \(H\).  Atomic formulas hold by the
substructure relation, and the Boolean steps are immediate.  For
\(\exists y\,\chi(\bar h,y)\), a witness in \(H\) is a witness in \(A\).
Conversely, if \(\exists y\,\chi(\bar h,y)\) holds in \((\mathfrak A,q)\),
the Skolem term \(t\) for \(\chi\) gives the witness
\(t^{\mathfrak A}(\bar h)\in H\).

For \(Qy\,\chi(\bar h,y)\), the induction hypothesis gives
\(\chi(H,\bar h)=\chi(A,\bar h)\cap H\).  If \(\chi(A,\bar h)\in q\), this set
is in \(q_H\).  Conversely, suppose \(\chi(A,\bar h)\cap H=\psi(A,\bar
c)\cap H\) with \(\bar c\) from \(H\) and \(\psi(A,\bar c)\in q\), but
\(\chi(A,\bar h)\notin q\).  Then \((A\setminus\chi(A,\bar h))\cap\psi(A,\bar
c)\in q\), so \(\exists y\,(\neg\chi(\bar h,y)\wedge\psi(\bar c,y))\) holds by
\cref{thm:p03-w03-p03-030-ultrafilter-facts}.  The Skolem term for this
formula gives a witness in \(H\) that lies in \(\psi(A,\bar c)\cap H\) and
not in \(\chi(A,\bar h)\cap H\), contradicting the equality.

So \((\mathfrak H,q_H)\preccurlyeq(\mathfrak A,q)\), and in particular
\((\mathfrak H,q_H)\entails T\).  Each member \(\phi(A,\bar h)\cap H\) of
\(q_H\) equals \(\phi(H,\bar h)\), so \((\mathfrak H,q_H)\) is reduced.
\end{proof}
\end{annexitem}

\begin{annexitem}{annex:p03-w03-p03-030-chains}{thm:p03-w03-p03-030-chains}{Unions of elementary chains of reduced structures}
\begin{proof}
For a formula \(\phi(\bar y)\) and \(\bar b\) from \(A\), write
\(\entails^*\phi(\bar b)\) when \((\mathfrak A_\nu,q_\nu)\entails\phi(\bar
b)\) for some \(\nu\) with \(\bar b\) in \(A_\nu\).  By elementarity of the
chain, this holds for one such \(\nu\) iff it holds for all of them.  Put
\(X_{\phi,\bar b}=\{a\in A:\ \entails^*\phi(a,\bar b)\}\).  The family \(q\)
of the statement is \(\{X_{\phi,\bar b}:\ \entails^*Qx\,\phi(x,\bar b)\}\).

We show by induction on \(\phi\) that \((\mathfrak A,q)\entails\phi(\bar b)\)
iff \(\entails^*\phi(\bar b)\).  Atomic formulas are decided in any
\(\mathfrak A_\nu\) containing \(\bar b\), because \(\mathfrak A\) is the union
of the \(L\)-structures.  Boolean steps are immediate.  For \(\exists
y\,\phi(\bar b,y)\): a witness \(a\) lies in some \(A_\mu\) with \(\mu\) at
least the stage of \(\bar b\), and then \((\mathfrak A_\mu,q_\mu)\entails
\exists y\,\phi(\bar b,y)\).  Conversely, a witness in the stage containing
\(\bar b\) is a witness in \(A\).

For \(Qy\,\phi(\bar b,y)\): by the induction hypothesis the set defined by
\(\phi(\bar b,y)\) in \((\mathfrak A,q)\) is \(X_{\phi,\bar b}\).  If
\(\entails^*Qy\,\phi(\bar b,y)\), then \(X_{\phi,\bar b}\in q\) by
definition.  Conversely, let \(X_{\phi,\bar b}=X_{\psi,\bar c}\) with
\(\entails^*Qy\,\psi(\bar c,y)\).  Take \(\nu\) with \(\bar b,\bar c\) in
\(A_\nu\).  If \((\mathfrak A_\nu,q_\nu)\entails\neg Qy\,\phi(\bar b,y)\),
then, as in the proof of \cref{thm:p03-w03-p03-030-hull},
\((\mathfrak A_\nu,q_\nu)\entails\exists y\,(\neg\phi(\bar
b,y)\wedge\psi(\bar c,y))\).  A witness \(a\in A_\nu\) then lies in
\(X_{\psi,\bar c}\setminus X_{\phi,\bar b}\), a contradiction.  Hence
\(\entails^*Qy\,\phi(\bar b,y)\).

Thus \((\mathfrak A_\nu,q_\nu)\preccurlyeq(\mathfrak A,q)\) for every
\(\nu\).  Each member \(X_{\phi,\bar b}\) of \(q\) is the set defined by
\(\phi(x,\bar b)\) in \((\mathfrak A,q)\), so \((\mathfrak A,q)\) is reduced.

For the second statement, argue by induction on \(\nu\le\delta\) that
\((\mathfrak A_\mu,q_\mu)\preccurlyeq(\mathfrak A_\nu,q_\nu)\) for all
\(\mu\le\nu\) and that \((\mathfrak A_\nu,q_\nu)\) is a reduced model of the
theory.  At a successor stage use parts (3) and (4) of
\cref{thm:p03-w03-p03-030-canonical-extension} and transitivity of
\(\preccurlyeq\).  At a limit stage the earlier stages form an elementary
chain of reduced structures satisfying the ultrafilter schemes, and the
first part applies.  An elementary extension satisfies the same sentences.
\end{proof}
\end{annexitem}

\begin{annexitem}{annex:p03-w03-p03-030-end-extension}{thm:p03-w03-p03-030-end-extension}{End-extensions and least new elements}
\begin{proof}
(1) Let \(b=[g]\in B\) and \(a\in A\) with \(b<a\) in \(\mathfrak B\).  By
part (2) of \cref{thm:p03-w03-p03-030-canonical-extension},
\(\{x:g(x)<a\}\in q\).  Let \(\gamma(x,y,\bar p)\) define the graph of \(g\),
with parameters \(\bar p\).  Then \((\mathfrak A,q)\entails Qx\,\exists
y\,(\gamma(x,y,\bar p)\wedge y<a)\).  The instance of (B) for
\(\phi=\gamma\), with \(z\defeq a\), gives \(c\in A\) with
\((\mathfrak A,q)\entails Qx\,\gamma(x,c,\bar p)\), that is,
\(\{x:g(x)=c\}\in q\).  So \(b=[c_c]\), and \(b\in A\) under the
identification.

(2) First, (B) holds.  Let \((\mathfrak A,q)\) be a model of \(T\) and
suppose \(Qx\,\exists y\,(\phi(x,y)\wedge y<z)\) for some \(z\) (with
parameters).  Since \(\{x:z<x\}\in q\), the intersection gives
\(Qx\,\exists y\,(\phi(x,y)\wedge y<z\wedge z<x)\).  By transitivity of
\(<\) and supersets, \(Qx\,\exists y\,(\phi(x,y)\wedge y<x)\) holds.  (R) then
gives \(\exists y\,Qx\,\phi(x,y)\).

Now let \((\mathfrak B,r)\) be the canonical extension.  For \(a\in A\),
\(\{x:a<x\}\in q\), so \(a<\iota\) by the \L o\'s equivalence; and
\(\iota\notin A\) by \cref{thm:p03-w03-p03-030-large-sets}.  Let \(b=[g]\in
B\setminus A\) and suppose \(b<\iota\).  Then \(\{x:g(x)<x\}\in q\), so with
\(\gamma\) as in (1), \(Qx\,\exists y\,(\gamma(x,y,\bar p)\wedge y<x)\).  By (R)
there is \(c\) with \(\{x:g(x)=c\}\in q\), so \(b=c\in A\), a contradiction.
Since \(<\) is linear in \(\mathfrak B\), every element of \(B\setminus A\) is
at least \(\iota\).  So \(\iota\) is the least element of \(B\setminus A\).
\end{proof}
\end{annexitem}

\begin{annexitem}{annex:p03-w03-p03-030-strongly-like}{thm:p03-w03-p03-030-strongly-like}{Strongly \(\lambda\)-like models from a self-extending theory}
\begin{proof}
By \cref{thm:p03-w03-p03-030-chains} each \((\mathfrak A_\nu,q_\nu)\),
\(\nu\le\lambda\), is a reduced model of \(T'\), and the chain is elementary.
By \cref{thm:p03-w03-p03-030-end-extension}(2), \(T'\) proves (B), and the
least new element \(a_\nu\) exists.

\emph{Initial segments.}  We show by induction on \(\nu\le\lambda\) that
\(A_\mu\) is an initial segment of \(A_\nu\) for all \(\mu\le\nu\).  At
successors this follows from \cref{thm:p03-w03-p03-030-end-extension}(1) and
transitivity.  At a limit \(\nu\), let \(c\in A_\nu\) with \(c<a\) for some
\(a\in A_\mu\), \(\mu<\nu\).  Then \(c\in A_\rho\) for some \(\rho<\nu\).
Both \(c\) and \(a\) lie in \(A_{\max(\mu,\rho)}\), in which \(A_\mu\) is an
initial segment, so \(c\in A_\mu\).

\emph{Sizes.}  \(\abs{A_0}<\lambda\).  If \(\abs{A_\nu}<\lambda\), then
\(\abs{A_{\nu+1}}\le\abs{A_\nu}+\abs L+\aleph_0<\lambda\) by part (5) of
\cref{thm:p03-w03-p03-030-canonical-extension}, since
\(\lambda>\abs L+\aleph_0\) and \(\lambda\) is an infinite cardinal.  At a
limit \(\nu<\lambda\), \(A_\nu\) is the union of \(\abs\nu<\lambda\) sets,
each of size less than \(\lambda\).  Regularity of \(\lambda\) makes the
supremum \(\kappa'\) of their sizes less than \(\lambda\), so
\(\abs{A_\nu}\le\abs\nu\cdot\kappa'+\aleph_0<\lambda\).  The elements \(a_\nu\)
are distinct, since \(a_\nu\notin A_\nu\supseteq\{a_\mu:\mu<\nu\}\).  So
\(\abs{A_\lambda}\ge\lambda\), and \(\abs{A_\lambda}\le\lambda\cdot\lambda=
\lambda\).

(1) If \(x\in A_\nu\) then \(x<a_\nu\): \(a_\nu\notin A_\nu\), \(A_\nu\) is an
initial segment of \(A_\lambda\), and \(<\) is linear.  If \(x\in A_\lambda\)
and \(x<a_\nu\), then \(x\in A_{\nu+1}\), because \(A_{\nu+1}\) is an initial
segment of \(A_\lambda\) containing \(a_\nu\).  If \(x\) were not in
\(A_\nu\), it would be an element of \(A_{\nu+1}\setminus A_\nu\) below its
least element.  So \(x\in A_\nu\).

(2) If \(\mu<\nu\), then \(a_\mu\in A_{\mu+1}\subseteq A_\nu\), so
\(a_\mu<a_\nu\) by (1).  Let \(\alpha<\lambda\) be a limit.  The element
\(a_\alpha\) is an upper bound of \(\{a_\nu:\nu<\alpha\}\).  If \(c\) is an
upper bound with \(c<a_\alpha\), then \(c\in A_\alpha\) by (1), so
\(c\in A_\nu\) for some \(\nu<\alpha\), and \(c<a_\nu\) by (1), a
contradiction.  So \(a_\alpha\) is the least upper bound.

(3) The \(L\)-reduct of \(\mathfrak A_\lambda\) satisfies every first-order
sentence of \(T'\), because \((\mathfrak A_0,q_0)\preccurlyeq(\mathfrak
A_\lambda,q_\lambda)\).  It has cardinality \(\lambda\).  For \(c\in
A_\lambda\), choose \(\nu<\lambda\) with \(c\in A_\nu\).  Then
\(\{x:x<c\}\subseteq A_\nu\) because \(A_\nu\) is an initial segment, so this
set has size less than \(\lambda\).  Hence the order is \(\lambda\)-like.
The map \(e(\nu)=a_\nu\) is increasing on \(\lambda\), and by (2) it satisfies
the continuity condition of \cref{def:p03-w03-p03-030-cardinal-like} at
every nonzero \(\alpha=\sup\{\nu:\nu<\alpha\}\), that is, at every limit
\(\alpha\).  So the order is \((\lambda,\lambda)\)-like, which is strongly
\(\lambda\)-like.  Its range is also cofinal: every \(c\in A_\nu\) lies below
\(a_\nu\).  It is closed in the sense of (2).
\end{proof}
\end{annexitem}

\begin{annexitem}{annex:p03-w03-p03-030-indiscernibles}{thm:p03-w03-p03-030-indiscernibles}{The new elements form an indiscernible sequence of any prescribed length}
\begin{proof}
By \cref{thm:p03-w03-p03-030-chains},
\((\mathfrak A_\mu,q_\mu)\preccurlyeq(\mathfrak A_\nu,q_\nu)\) whenever
\(\mu\le\nu\le\delta\).  We prove the displayed equivalence by induction on
\(n\), for all formulas, all \(\delta\) and all increasing
\(\nu_1<\cdots<\nu_n<\delta\) simultaneously.

For \(n=0\) the claim is \((\mathfrak A_\delta,q_\delta)\entails\phi(\bar
b)\iff(\mathfrak A_0,q_0)\entails\phi(\bar b)\), which is elementarity.

Let \(n\ge1\).  The elements \(a_{\nu_1},\ldots,a_{\nu_{n-1}}\) lie in
\(A_{\nu_n}\), since \(\nu_i+1\le\nu_n\) for \(i<n\).  The element \(a_{\nu_n}\)
is the class of the identity function in the canonical extension
\((\mathfrak A_{\nu_n+1},q_{\nu_n+1})\) of \((\mathfrak A_{\nu_n},q_{\nu_n})\).
By elementarity and part (2) of
\cref{thm:p03-w03-p03-030-canonical-extension}, applied with constant
functions for the other arguments and the identity for the last one,
\[
 (\mathfrak A_\delta,q_\delta)\entails\phi(a_{\nu_1},\ldots,a_{\nu_n},\bar b)
 \iff
 (\mathfrak A_{\nu_n},q_{\nu_n})\entails
   Qx_n\,\phi(a_{\nu_1},\ldots,a_{\nu_{n-1}},x_n,\bar b).
\]
The chain \((\mathfrak A_\nu,q_\nu)_{\nu\le\nu_n}\) is of the same kind, with
top stage \(\nu_n\).  The induction hypothesis for \(n-1\), applied to the
formula \(Qx_n\,\phi\), turns the right side into \((\mathfrak A_0,q_0)\entails
Qx_1\cdots Qx_{n-1}\,Qx_n\,\phi(x_1,\ldots,x_n,\bar b)\).  When \(n=1\) the
case \(n-1=0\) is elementarity of the stage \(\nu_1\) over stage \(0\).

The right side of the equivalence does not depend on the indices, so any two
increasing \(n\)-tuples from the sequence satisfy the same formulas with
parameters from \(A_0\).  This is order indiscernibility over \(A_0\).  The
\(a_\nu\) are distinct, because \(a_\nu\notin A_\nu\) while \(a_\mu\in A_\nu\)
for \(\mu<\nu\).
\end{proof}
\end{annexitem}

\begin{annexitem}{annex:p03-w03-p03-030-definable-ultrapower-n}{thm:p03-w03-p03-030-definable-ultrapower-n}{A self-extending expansion of the natural numbers}
\begin{proof}
The interpretations of the symbols of \(L_{k+1}\) are fixed by recursion on
\(k\).  The formula \(\chi\) belongs to \(L_k(Q)\), and its truth in
\((\mathfrak N\restriction L_k,\mathcal U)\) is already determined.  Every
nonempty subset of \(\mathbb N\) has a least element, so \(f_\chi\) is well
defined.  Because \(Q\) only asks about definable sets, the structures
\((\mathfrak N,\mathcal U)\) and \((\mathfrak N,q)\) satisfy the same
formulas.

(1) Every \(L_\omega(Q)\)-formula \(\chi(\bar z,y)\) lies in some
\(L_k(Q)\).  Then \(\forall\bar z\,(\exists y\,\chi\rightarrow\chi(\bar
z,f_\chi(\bar z)))\) holds, since \(f_\chi(\bar z)\) is the least witness
whenever one exists.  So \(T\) is a Skolem theory.  \(\mathcal U\) is an
ultrafilter on all subsets of \(\mathbb N\), so the complement, superset and
intersection schemes hold.  It contains no singleton, because it is
nonprincipal.  Hence \(T\) is self-extending; it is consistent, since
\((\mathfrak N,q)\) is a model.  A finite set \(\{n_1,\ldots,n_k\}\) is not in
\(\mathcal U\): otherwise, by the ultrafilter property, one of the
singletons would be.  So every cofinite set, in particular
\(\{x:x>n\}\), is in \(\mathcal U\), and \(\forall y\,Qx\,(y<x)\in T\).

(2) Suppose \([g]<a\) with \(a\in\mathbb N\).  Then \(X=\{x:g(x)<a\}\in q\).
\(X\) is the union of the finitely many definable sets
\(X_c=\{x:g(x)=c\}\), \(c<a\).  If none of them were in \(q\), their
complements would all be, and so would the intersection of these
complements.  That intersection is disjoint from \(X\), which contradicts
\(\varnothing\notin q\).  So some \(X_c\in q\), and \([g]=c\).

(3) Let \(\phi(x,y)\) be \(y+y=x\vee y+y+1=x\).  For \(x\ge1\), the only
\(y\) with \(\phi(x,y)\) is \(\lfloor x/2\rfloor<x\).  So the cofinite set
\(\{x:\exists y\,(\phi(x,y)\wedge y<x)\}\) is in \(q\).  For each \(c\),
\(\{x:\phi(x,c)\}=\{2c,2c+1\}\) is finite, hence not in \(q\).  So the instance
of (R) for \(\phi\) fails.

Let \(b=[g]\) be new, that is, \(b\notin\mathbb N\).  Then
\(\{x:g(x)\ge1\}\in q\), for otherwise \(\{x:g(x)=0\}\in q\) and \(b=0\).
Put \(h(x)=\lfloor g(x)/2\rfloor\), a definable function.  On
\(\{x:g(x)\ge1\}\) we have \(h(x)<g(x)\), so \([h]<b\).  If \(\{x:h(x)=c\}\in
q\) for some \(c\), then \(\{x:g(x)=2c\}\cup\{x:g(x)=2c+1\}\in q\), so one of
these two sets is in \(q\) and \(b\in\mathbb N\), a contradiction.  So \([h]\)
is new and smaller than \(b\).  Hence no new element is least.  New elements
exist, since \(\iota\) is one by \cref{thm:p03-w03-p03-030-large-sets}.
\end{proof}
\end{annexitem}

\begin{annexitem}{annex:p03-w03-p03-030-tuple-orders}{thm:p03-w03-p03-030-tuple-orders}{Arithmetic of the tuple orders}
\begin{proof}
We first record the case \(n=1\).
\begin{enumerate}
\item[(a)] \emph{A nonempty definable set \(D\subseteq M\) with an upper bound
  \(b\) has a greatest element.}  Otherwise every element of \(D\) has a
  larger element in \(D\).  The induction scheme then makes \(D\) cofinal,
  which gives an element of \(D\) above \(b\), a contradiction.
\item[(b)] \emph{Every \(a\) has an immediate successor \(a^+\).}  The set
  \(\{x:a<x\}\) is nonempty, since there is no last element; its least element
  under the least-number principle is \(a^+\).
\item[(c)] \emph{Every \(a\ne0\) has an immediate predecessor.}  The set
  \(\{x:x<a\}\) contains \(0\) and is bounded by \(a\); its greatest element,
  by (a), is the predecessor.
\end{enumerate}
Now the parts of the lemma, by induction on \(n\); for \(n=1\) parts (1),
(2) and (5) are immediate, and (3) and (4) are (a)--(c).

(1) The map \(\bar x\mapsto(\max\bar x,\,x_n,\,(x_0,\ldots,x_{n-1}))\) is
injective, and \(<_{n+1}\) is the lexicographic order of its values under
\(<\), \(<\) and \(<_n\).  The lexicographic product of strict linear orders
is a strict linear order.  \(0^{n+1}\) is least: a tuple with maximum \(0\)
is \(0^{n+1}\), and every other tuple has a larger maximum.

(2) Let \(S\neq\varnothing\) be a definable set of \((n+1)\)-tuples.  Let
\(m\) be least in \(\{\max\bar x:\bar x\in S\}\), let \(c\) be least in
\(\{x_n:\bar x\in S,\ \max\bar x=m\}\), and let \(\bar u\) be \(<_n\)-least in
\(S'=\{(x_0,\ldots,x_{n-1}):(x_0,\ldots,x_{n-1},c)\in S,\
\max(x_0,\ldots,x_{n-1},c)=m\}\), by the induction hypothesis.  All three sets
are definable and nonempty.  Comparing the three keys in turn shows that
\((\bar u,c)\) is \(<_{n+1}\)-least in \(S\).

(3) Let \(S\neq\varnothing\) be definable with a \(<_{n+1}\)-upper bound
\(\bar w\).  For \(\bar x\in S\) we have \(\max\bar x\le\max\bar w\), so the
set of maxima has an upper bound and a greatest element \(m\), by (a).  The
set \(\{x_n:\bar x\in S,\ \max\bar x=m\}\) is bounded by \(m\) and has a
greatest element \(c\).  The set \(S'\), formed with this \(c\), consists of
tuples all of whose entries are at most \(m\).  By (5) for \(n\), \(S'\) is
\(<_n\)-bounded by \((m^+)^n\), so it has a greatest element \(\bar u\) by the
induction hypothesis.  Then \((\bar u,c)\) is greatest in \(S\).

(4) For \(\bar u\), the set \(\{\bar x:\bar u<_{n+1}\bar x\}\) contains
\(((\max\bar u)^+)^{n+1}\) by (5), so it has a least element by (2), which is
the immediate successor.  For \(\bar u\ne0^{n+1}\), the set
\(\{\bar x:\bar x<_{n+1}\bar u\}\) contains \(0^{n+1}\) and is bounded by
\(\bar u\), so it has a greatest element by (3), which is the immediate
predecessor.

(5) If \(\max\bar x<z=\max z^{n+1}\), the first disjunct of the definition
gives \(\bar x<_{n+1}z^{n+1}\).  If \(z<z'\), the same disjunct gives
\(z^{n+1}<_{n+1}z'^{n+1}\).

(6) Let \(D=\{u:\exists b\,\forall z\,(z\le u\rightarrow F(z)<b)\}\), a
definable set.  \(0\in D\), with \(b=F(0)^+\).  If \(u\in D\) with bound \(b\),
then \(u^+\in D\): the elements \(z\le u^+\) are those with \(z\le u\) and
\(u^+\) itself, so \(\max(b,F(u^+))^+\) is a bound.  \(D\) is closed
downwards.  Suppose \(m\notin D\).  Then every element of \(D\) is below \(m\),
since otherwise \(m\in D\) by downward closure.  So \(D\) is nonempty and
bounded, and has a greatest element \(d\) by (a).  But \(d^+\in D\), a
contradiction.  So \(m\in D\).
\end{proof}
\end{annexitem}

\begin{annexitem}{annex:p03-w03-p03-030-tree}{thm:p03-w03-p03-030-tree}{The leftmost large branch}
\begin{proof}
Ranks are \((n+1)\)-tuples ordered by \(<_{n+1}\), and
\cref{thm:p03-w03-p03-030-tuple-orders} is used for them without further
mention.  Write \(C(\bar u,y)=\{z:E(\bar u,y,z)\}\).  A definable set \(X\)
is \emph{large} when \(\forall w\,\exists z>w\,(\theta(z)\wedge z\in X)\).
Thus the first conjunct of \(S(\bar u,y)\) says that \(C(\bar u,y)\) is
large, and the second says that every \(v\) with \(L(\bar u,v,y)\) has
\(C(\bar u,v)\) not large.

\emph{Basic facts.}
\begin{enumerate}
\item[(i)] For each \(\bar u\), \(E(\bar u,\cdot,\cdot)\) is an equivalence
  relation.  If \(\bar u\le_{n+1}\bar u'\), then \(E(\bar u',y,z)\) implies
  \(E(\bar u,y,z)\), since fewer parameter tuples are compared.
\item[(ii)] Supersets of large sets are large.  If \(X_1\) and \(X_2\) are not
  large, witnessed by \(w_1\) and \(w_2\), then \(\max(w_1,w_2)\) witnesses that
  \(X_1\cup X_2\) is not large.
\item[(iii)] \(L(\bar u,y,z)\) depends only on the classes \(C(\bar u,y)\) and
  \(C(\bar u,z)\).  A witness \(\bar w<_{n+1}\bar u\) sees the same values of
  \(\phi(\bar w,\cdot)\), and of \(E(\bar w,\cdot,\cdot)\), on
  \(E(\bar u)\)-equivalent elements.  The same holds for \(S(\bar u,y)\),
  which is built from \(L\) and \(C\).
\item[(iv)] \emph{Splitting.}  Let \(\bar u'\) be the immediate successor of
  \(\bar u\).  Then \(\{\bar w:\bar w<_{n+1}\bar u'\}=\{\bar w:\bar
  w\le_{n+1}\bar u\}\), so \(C(\bar u',y)=C(\bar u,y)\cap\{z:\phi(\bar
  u,z)\leftrightarrow\phi(\bar u,y)\}\).  Hence \(C(\bar u,y)\) is the union
  of \(C(\bar u',y)\) and the set of its members \(z\) with
  \(\phi(\bar u,z)\not\leftrightarrow\phi(\bar u,y)\).  The latter, if
  nonempty, is \(C(\bar u',y'')\) for any of its members \(y''\).
\end{enumerate}

\emph{Proof of (1).}  \(L(\bar u,y,y)\) is impossible.  Let \(y,z\) be
inequivalent at \(\bar u\).  The set of \(\bar w<_{n+1}\bar u\) with
\(\phi(\bar w,y)\not\leftrightarrow\phi(\bar w,z)\) is nonempty and definable,
so it has a least element \(\bar w_0\).  Then \(E(\bar w_0,y,z)\), and
\(L(\bar u,y,z)\) or \(L(\bar u,z,y)\) holds, depending on which of \(y,z\)
satisfies \(\phi(\bar w_0,\cdot)\).  Any witness to \(L(\bar u,y,z)\) is a
point of difference below which \(y\) and \(z\) agree, so it is \(\bar
w_0\).  Hence \(L(\bar u,y,z)\) and \(L(\bar u,z,y)\) cannot both hold.  If
\(y,z\) are equivalent at \(\bar u\), neither holds.

\emph{Transitivity.}  Let \(L(\bar u,y,z)\) with witness \(\bar w_1\) and
\(L(\bar u,z,v)\) with witness \(\bar w_2\).  Since \(\phi(\bar w_1,z)\) and
\(\neg\phi(\bar w_2,z)\), we have \(\bar w_1\ne\bar w_2\).  If \(\bar
w_1<_{n+1}\bar w_2\), then \(E(\bar w_1,z,v)\) by (i), so \(E(\bar
w_1,y,v)\).  Also \(\phi(\bar w_1,v)\leftrightarrow\phi(\bar w_1,z)\), which
is true, while \(\neg\phi(\bar w_1,y)\).  So \(\bar w_1\) witnesses \(L(\bar
u,y,v)\).  If \(\bar w_2<_{n+1}\bar w_1\), then \(E(\bar w_2,y,z)\), so
\(E(\bar w_2,y,v)\).  Also \(\phi(\bar w_2,y)\leftrightarrow\phi(\bar
w_2,z)\), which is false, while \(\phi(\bar w_2,v)\).  So \(\bar w_2\)
witnesses \(L(\bar u,y,v)\).  With (iii), \(L(\bar u,\cdot,\cdot)\) is a strict
linear order of the classes.

\emph{Persistence.}  \emph{If \(\bar w\le_{n+1}\bar u\) and \(C(\bar w,y)\) is
large, then \(C(\bar w,y)\) contains a large class \(C(\bar u,y')\).}  Let
\(D\) be the definable set of ranks \(\bar u\ge_{n+1}\bar w\) for which some
\(y'\in C(\bar w,y)\) has \(C(\bar u,y')\) large.  If some rank
\(\ge_{n+1}\bar w\) is not in \(D\), let \(\bar u_0\) be the least one.
Since \(\bar w\in D\), \(\bar u_0>_{n+1}\bar w\), and its immediate
predecessor \(\bar u'\) satisfies \(\bar u'\ge_{n+1}\bar w\), so
\(\bar u'\in D\) by leastness.  Take \(y'\in C(\bar w,y)\) with
\(C(\bar u',y')\) large.  By (iv), \(C(\bar u',y')\) is a union of at most two
classes of rank \(\bar u_0\), both inside \(C(\bar w,y)\) by (i).  By (ii)
one of them is large.  So \(\bar u_0\in D\), a contradiction.

\emph{Proof of (2), existence.}  At the least rank \(0^{n+1}\) no tuple is
compared.  There is one class, all of \(M\); it is large because
\(\theta(M)\) is unbounded, and \(L(0^{n+1},\cdot,\cdot)\) is empty.  So
\(S(0^{n+1},y)\) holds for every \(y\).  Suppose some rank has no \(y\) with
\(S\), and let \(\bar u_0\) be the least such rank, with immediate
predecessor \(\bar u'\) and some \(y\) with \(S(\bar u',y)\).  Put
\(C=C(\bar u',y)\), which is large.  Put
\(C^0=\{z\in C:\neg\phi(\bar u',z)\}\) and
\(C^1=\{z\in C:\phi(\bar u',z)\}\).  By (iv) each of them is empty or a class
of rank \(\bar u_0\), and by (ii) one is large.  Let \(C^*\) be \(C^0\) if
\(C^0\) is large and \(C^1\) otherwise, and pick \(y^*\in C^*\).  Let \(v\)
satisfy \(L(\bar u_0,v,y^*)\), with witness \(\bar w\le_{n+1}\bar u'\).  If
\(\bar w<_{n+1}\bar u'\), then \(\bar w\) also witnesses \(L(\bar u',v,y^*)\).
Since \(C(\bar u',y^*)=C\) and \(S(\bar u',y)\) holds, \(C(\bar u',v)\) is not
large, and neither is its subset \(C(\bar u_0,v)\).  If \(\bar w=\bar u'\),
then \(v\in C\), \(\neg\phi(\bar u',v)\) and \(\phi(\bar u',y^*)\).  So
\(C(\bar u_0,v)=C^0\) and \(C^*=C^1\), which by the choice of \(C^*\) means that
\(C^0\) is not large.  In both cases \(C(\bar u_0,v)\) is not large, so
\(S(\bar u_0,y^*)\) holds, a contradiction.

\emph{Proof of (2), uniqueness.}  If \(S(\bar u,y)\) and \(S(\bar u,y')\) with
\(y,y'\) inequivalent, then by (1), say, \(L(\bar u,y,y')\).  The second
conjunct of \(S(\bar u,y')\) makes \(C(\bar u,y)\) not large, contrary to the
first conjunct of \(S(\bar u,y)\).  By (iii) the set defined by \(S(\bar
u,\cdot)\) is a union of classes, so it is exactly one class, and that class
is large.

\emph{Proof of (3).}  Let \(S(\bar u,y)\) and \(\bar w<_{n+1}\bar u\).  The
class \(C(\bar w,y)\) contains the large class \(C(\bar u,y)\), so it is
large.  Suppose \(L(\bar w,v,y)\) with \(C(\bar w,v)\) large.  By
persistence, \(C(\bar w,v)\) contains a large class \(C(\bar u,v')\).  By
(iii), \(L(\bar w,v',y)\) holds, with a witness \(\hat w<_{n+1}\bar w<_{n+1}\bar
u\), which also witnesses \(L(\bar u,v',y)\).  Then \(S(\bar u,y)\) says that
\(C(\bar u,v')\) is not large, a contradiction.  So \(S(\bar w,y)\).

\emph{Proof of (4).}  Let \(\bar x<_{n+1}\bar u\) and \(S(\bar u,y)\).
Suppose \(\sigma(\bar x)\).  Taking \(\bar w=\bar u\) gives \(y'\) with
\(S(\bar u,y')\) and \(\phi(\bar x,y')\).  By (2), \(E(\bar u,y,y')\), and since
\(\bar x<_{n+1}\bar u\), \(\phi(\bar x,y)\) holds.  Conversely, suppose
\(\phi(\bar x,y)\) and let \(\bar w\) be any rank.  If \(\bar w\le_{n+1}\bar u\),
then \(S(\bar w,y)\) by (3), and \(y\) itself is the required witness.  If
\(\bar u<_{n+1}\bar w\), take \(y'\) with \(S(\bar w,y')\), by (2).  By (3),
\(S(\bar u,y')\), so \(E(\bar u,y,y')\) by (2), and therefore \(\phi(\bar
x,y')\).  Hence \(\sigma(\bar x)\).
\end{proof}
\end{annexitem}

\begin{annexitem}{annex:p03-w03-p03-030-mss}{thm:p03-w03-p03-030-mss}{Self-extending theories from the induction scheme}
\begin{proof}
Below, ``\(T\) proves'' a first-order sentence means that the sentence holds
in every model of \(T\).  By the completeness theorem, this is the same as
derivability.

\emph{Enumeration.}  The vocabulary is countable, so the first-order formulas
can be listed as \(\psi_0,\psi_1,\ldots\), each with free variables among
\(x_0,x_1,\ldots\) and \(y\).  Let \(k_j\) be the largest index of an \(x\)
free in \(\psi_j\) (or \(0\)).  Set \(n_{-1}=-1\), and put \(\psi_j\) at position
\(n_j=\max(k_j,n_{j-1}+1)\), and fill the other positions with \(y=y\).  The
result is a list \((\phi_n)_{n<\omega}\) in which \(\phi_n\) has its free
variables among \(x_0,\ldots,x_n,y\) and every formula occurs
\cite[p.~207]{Schmerl1985TransferTheorems}.

\emph{Recursion.}  We define first-order formulas \(\theta_n(y)\) and
\(\sigma_n(x_0,\ldots,x_n)\), without parameters, such that \(T\) proves
\begin{gather*}
 \text{(C1) }\forall w\,\exists y>w\,\theta_n(y),\qquad
 \text{(C2) }\forall y\,(\theta_n(y)\rightarrow\theta_{n-1}(y)),\\
 \text{(C3) }\forall\bar x\,\exists w\,\forall y>w\,
        (\theta_n(y)\rightarrow(\phi_n(\bar x,y)\leftrightarrow\sigma_n(\bar x))).
\end{gather*}
These are the three conditions of the source
\cite[p.~207]{Schmerl1985TransferTheorems}.  Let \(\theta_{-1}(y)\) be
\(y=y\); (C1) holds for it because there is no last element.  Given
\(\theta_{n-1}\) satisfying (C1), form \(E\), \(L\), \(S\) and \(\sigma_n\) as
in \cref{def:p03-w03-p03-030-tree-search}, with \(\phi=\phi_n\) and
\(\theta=\theta_{n-1}\).  Let
\(\mathrm{Ag}(z,y)\defeq\forall\bar x<_{n+1}z^{n+1}\,(\sigma_n(\bar
x)\leftrightarrow\phi_n(\bar x,y))\), and
\begin{multline*}
 \theta_n(y)\defeq\theta_{n-1}(y)\wedge\exists z<y\,\bigl[\mathrm{Ag}(z,y)\\
  \wedge\forall w\,\bigl((\mathrm{Ag}(z,w)\wedge\theta_{n-1}(w))\rightarrow
  w\le z\vee y\le w\bigr)\bigr].
\end{multline*}
This is the source's formula with its last disjunct read as \(y\le w\); see
\cref{def:p03-w03-p03-030-limits}.

Fix a model \(M\) of \(T\).  \cref{thm:p03-w03-p03-030-tree} applies, since
\(\theta_{n-1}(M)\) is unbounded.  For \(z\in M\), take \(y_0\) with
\(S(z^{n+1},y_0)\).  Its class is large, so it contains some \(w>z\) with
\(\theta_{n-1}(w)\).  Then \(S(z^{n+1},w)\) holds, and by part (4) of the
lemma, \(\mathrm{Ag}(z,w)\).  Let \(F(z)\) be the least \(w>z\) with
\(\theta_{n-1}(w)\wedge\mathrm{Ag}(z,w)\).  Then \(F\) is a definable total
function, and
\[
 \theta_n(y)\iff\exists z<y\ \ y=F(z).
\]
If \(y=F(z)\), then \(z\) is a witness: any \(w\) with \(\mathrm{Ag}(z,w)\wedge
\theta_{n-1}(w)\) and \(w>z\) satisfies \(w\ge F(z)\).  Conversely, if
\(z<y\) witnesses \(\theta_n(y)\), then \(F(z)\le y\), and the clause for
\(w=F(z)\) gives \(F(z)\le z\), which is false, or \(y\le F(z)\).

(C1): \(F(w)>w\) and \(\theta_n(F(w))\).  (C2) is built into \(\theta_n\).
(C3): given \(\bar x\), let \(m=\max\bar x\).  By part (6) of
\cref{thm:p03-w03-p03-030-tuple-orders} there is \(b\) with \(F(z)<b\) for all
\(z\le m\).  If \(y>b\) and \(\theta_n(y)\), then \(y=F(z)\) for some \(z\),
and \(z>m\).  So \(\bar x<_{n+1}z^{n+1}\) by part (5), and \(\mathrm{Ag}(z,y)\)
gives \(\phi_n(\bar x,y)\leftrightarrow\sigma_n(\bar x)\).

By (C2), \(T\) proves \(\theta_k\rightarrow\theta_j\) whenever \(j\le k\).

\emph{The formulas \(\sigma_\phi\).}  Let \(\phi(x_0,\ldots,x_{k-1},y)\) be a
first-order formula, and let \(n\) be a position with \(\phi_n=\phi\).  If
\(\bar b,\bar b'\) agree on the first \(k\) coordinates, then in any \(M\)
we have \(\sigma_n(\bar b)\leftrightarrow\sigma_n(\bar b')\).  To see this,
choose \(y\) above the two bounds of (C3) with \(\theta_n(y)\), by (C1); then
\(\sigma_n(\bar b)\leftrightarrow\phi(\bar b,y)\leftrightarrow\sigma_n(\bar
b')\).  Put \(\sigma_\phi(x_0,\ldots,x_{k-1})\defeq\forall
u_k\cdots\forall u_n\,\sigma_n(x_0,\ldots,x_{k-1},u_k,\ldots,u_n)\).  This is
equivalent in every model of \(T\) to \(\sigma_n\) with any values in the
padding positions.

\emph{Key property.}  In every \(M\models T\), for every first-order
\(\phi(\bar x,y)\) and \(\bar b\): \(M\entails\sigma_\phi(\bar b)\) iff there
are \(j\) and \(w\) with \(\forall y>w\,(\theta_j(y)\rightarrow\phi(\bar
b,y))\).  If \(\sigma_\phi(\bar b)\), take \(j=n\) and use (C3).  Conversely,
suppose such \(j,w\) exist but \(\neg\sigma_\phi(\bar b)\).  By (C3) there is
\(w'\) with \(\forall y>w'\,(\theta_n(y)\rightarrow\neg\phi(\bar b,y))\).  By
(C1) there is \(y>\max(w,w')\) with \(\theta_{\max(j,n)}(y)\).  Then both
\(\phi(\bar b,y)\) and \(\neg\phi(\bar b,y)\), a contradiction.

\emph{The family \(q_M\).}  Let \(q_M\) be the set of first-order definable
\(X\subseteq M\) for which some \(j,w\) satisfy \(\forall y>w\,(\theta_j(y)
\rightarrow y\in X)\).  By the key property, \(\phi(M,\bar b)\in q_M\) iff
\(M\entails\sigma_\phi(\bar b)\).  The family \(q_M\) is closed under
definable supersets.  It is closed under intersections, using
\(\theta_{\max(j,j')}\) and \(\max(w,w')\).  It contains every final segment
\(\{x:x>b\}\), and it does not contain \(\varnothing\), by (C1).  For a
definable \(X=\phi(M,\bar b)\), either \(\sigma_\phi(\bar b)\) holds and
\(X\in q_M\), or it fails and (C3) puts \(M\setminus X\) in \(q_M\).  No
singleton \(\{b\}\) is in \(q_M\), since (C1) gives \(y>\max(w,b)\) with
\(\theta_j(y)\).

\emph{Elimination of \(Q\).}  Every \(\mathscr L(Q)\)-formula \(\chi(\bar z)\)
has a first-order \(\chi^*(\bar z)\), not depending on \(M\), with
\((M,q_M)\entails\forall\bar z\,(\chi\leftrightarrow\chi^*)\) for every
\(M\models T\).  The proof is by induction on \(\chi\).  Atomic formulas are
their own translations, and Boolean connectives and \(\exists\) commute with
the translation.  For \(Qy\,\chi(\bar z,y)\): \((M,q_M)\entails
Qy\,\chi(\bar c,y)\) iff \(\chi^*(M,\bar c)\in q_M\), iff
\(M\entails\sigma_{\chi^*}(\bar c)\).  So \((Qy\,\chi)^*\defeq\sigma_{\chi^*}\),
with variables renamed.  Consequently every set definable in \((M,q_M)\) is
first-order definable.  So \((M,q_M)\) is reduced, and \(q_M\) is a
nonprincipal ultrafilter on its definable sets containing all final
segments.

\emph{The theory \(T'\).}  Let \(T'\) be the set of \(\mathscr
L(Q)\)-sentences true in \((M,q_M)\) for every \(M\models T\).  It contains
\(T\), and it is consistent because \(T\) has a model.  Each \((M,q_M)\) is a
reduced model of \(T'\).  We check the required sentences in an arbitrary
\((M,q_M)\).

\emph{Ultrafilter schemes and \(\forall y\,Qx\,(y<x)\)} follow from the
properties of \(q_M\) just proved.

\emph{Skolem property.}  For an \(\mathscr L(Q)\)-formula \(\chi(\bar x,y)\),
let \(t\) be the \(L\)-term that the Skolem property of \(T\) provides for
\(\chi^*\).  Since \(\chi\leftrightarrow\chi^*\) holds in \((M,q_M)\), so does
\(\forall\bar x\,(\exists y\,\chi\rightarrow\chi(\bar x,t(\bar x)))\).

\emph{The sentences \(Qy\,\phi(\bar x,y)\leftrightarrow\sigma_\phi(\bar x)\)}
hold by the description of \(q_M\).

\emph{The bounded scheme (B).}  Let \(\phi(x,y,\bar p)\) be an \(\mathscr
L(Q)\)-formula, and let \(\bar p\) and \(z\) be from \(M\) with \((M,q_M)\entails
Qx\,\exists y\,(\phi(x,y,\bar p)\wedge y<z)\).  Replacing \(\phi\) by
\(\phi^*\) changes no set, so assume \(\phi\) is first-order.  Let
\(\psi(y,\bar p,x)\defeq\phi(x,y,\bar p)\), viewed as a formula with
parameter variables \(y,\bar p\) and last variable \(x\), and let \(m\) be
its position in the list.  Suppose, for a contradiction, that
\(\{x:\phi(x,c,\bar p)\}\notin q_M\) for every \(c<z\).  Then
\(\neg\sigma_\psi(c,\bar p)\) for each \(c<z\), and (C3) gives some \(w\) with
\(\forall x>w\,(\theta_m(x)\rightarrow\neg\phi(x,c,\bar p))\).  Let \(W(c)\)
be the least such \(w\) when \(c<z\), and \(0\) otherwise; \(W\) is a definable
function.  By part (6) of \cref{thm:p03-w03-p03-030-tuple-orders} there is
\(b\) with \(W(c)<b\) for all \(c\le z\).  The hypothesis gives \(j,w_0\) such
that every \(x>w_0\) with \(\theta_j(x)\) has some \(c<z\) with
\(\phi(x,c,\bar p)\).  Take \(x>\max(b,w_0)\) with
\(\theta_{\max(j,m)}(x)\), by (C1).  Then some \(c<z\) has \(\phi(x,c,\bar
p)\), but \(x>W(c)\) and \(\theta_m(x)\) give \(\neg\phi(x,c,\bar p)\).  So
some \(c<z\) has \(\{x:\phi(x,c,\bar p)\}\in q_M\), that is,
\((M,q_M)\entails\exists y\,Qx\,\phi(x,y,\bar p)\).

So \(T'\) is a self-extending theory containing all the displayed sentences,
and every model \(M\) of \(T\) expands to the reduced model \((M,q_M)\) of
\(T'\).
\end{proof}
\end{annexitem}

\begin{annexitem}{annex:p03-w03-p03-030-is-successor}{thm:p03-w03-p03-030-is-successor}{The least-number principle over the induction scheme}
\begin{proof}
Write \(y\le z\) for \(y<z\vee y=z\).  The induction scheme makes \(<\) a
linear order with first element \(0\) and no last element.

(1) Let \(D\) be nonempty, definable and bounded above by \(b\).  If \(D\) had
no greatest element, every element of \(D\) would have a larger element in
\(D\).  The instance of the induction scheme for the formula defining \(D\)
would then make \(D\) cofinal, giving \(d\in D\) with \(b<d\), a
contradiction.  Now let \(a\ne0\).  The set \(\{x:x<a\}\) is defined with
the parameter \(a\).  It contains \(0\), because \(0\le a\) and \(0\ne a\), and
it is bounded above by \(a\).  Let \(p\) be its greatest element.  Then
\(p<a\), and no \(x\) satisfies \(p<x<a\), because such an \(x\) would belong
to the set and exceed \(p\).  So \(p\) is the immediate predecessor of \(a\).

(2) Suppose the least-number principle holds, and let \(x\in M\).  The set
\(\{y:x<y\}\) is nonempty, since there is no last element.  Let \(y\) be its
least element.  If \(x<z\), then \(z<y\) is impossible, so \(y\le z\) by
linearity.  Thus \(y\) is an immediate successor of \(x\).

Conversely, suppose every element \(x\) has an immediate successor, written
\(x^+\).  Let \(\phi(x,\bar b)\) be satisfied by some element \(a\).  If
\(\phi(0,\bar b)\) holds, \(0\) is the least element satisfying \(\phi\),
since \(0\le y\) for every \(y\).  Otherwise let
\[
 D=\{z:\forall y\,(y\le z\rightarrow\neg\phi(y,\bar b))\}.
\]
The only \(y\le0\) is \(0\), so \(0\in D\).  If \(z\in D\), then \(a\le z\)
would give \(\neg\phi(a,\bar b)\); hence \(z<a\), and \(D\) is bounded above by
\(a\).  By (1), \(D\) has a greatest element \(m\).  Since \(m<m^+\),
\(m^+\notin D\), so some \(y\le m^+\) satisfies \(\phi(y,\bar b)\).  Because
\(m\in D\), \(y\le m\) is impossible, so \(m<y\le m^+\).  If \(y<m^+\), then
\(m<y<m^+\), which contradicts the choice of \(m^+\); hence \(y=m^+\) and
\(\phi(m^+,\bar b)\) holds.  Finally, if \(y'<m^+\), then \(y'\le m\),
because \(m<y'<m^+\) is impossible.  Since \(m\in D\), \(\neg\phi(y',\bar
b)\).  So \(m^+\) is the least element satisfying \(\phi\), and the
least-number principle holds.
\end{proof}
\end{annexitem}

\begin{annexitem}{annex:p03-w03-p03-030-end-extension-pa}{thm:p03-w03-p03-030-end-extension-pa}{Proper elementary end-extensions}
\begin{proof}
Let \(L\) be the countable vocabulary of \(M\).  Put \(L_0=L\).  Let
\(L_{k+1}\) add to \(L_k\) a function symbol \(f_\chi\) for each first-order
\(L_k\)-formula \(\chi(\bar z,y)\), interpreted in \(M\) as the least \(y\)
with \(\chi(\bar z,y)\) if there is one, and as \(0\) otherwise.  The
least-number principle supplies the least element, and the first element
\(0\) exists by the induction scheme.  Let \(L^{\mathrm{sk}}=\bigcup_kL_k\),
which is countable, and let \(M^{\mathrm{sk}}\) be the expansion.

Every \(L_{k+1}\)-formula is equivalent in \(M^{\mathrm{sk}}\) to an
\(L_k\)-formula.  The graph of \(f_\chi\) is defined by the \(L_k\)-formula
\[
 (\chi(\bar z,y)\wedge\forall y'<y\,\neg\chi(\bar z,y'))\vee
 (\neg\exists y'\,\chi(\bar z,y')\wedge\neg\exists u\,(u<y)).
\]
An atomic formula containing a term \(f_\chi(\bar s)\) is equivalent to
\(\exists v\,(\text{graph}(\bar s,v)\wedge\alpha')\), where \(\alpha'\) is the
atomic formula with that occurrence replaced by the new variable \(v\).
Repeating this for the occurrences of new symbols, innermost first, removes
all of them in finitely many steps.  Hence every \(L^{\mathrm{sk}}\)-formula
is equivalent to an \(L\)-formula, and \(M^{\mathrm{sk}}\) satisfies the
induction scheme and the least-number principle for all
\(L^{\mathrm{sk}}\)-formulas.  It is a Skolem structure: an
\(L^{\mathrm{sk}}\)-formula \(\chi\) lies in some \(L_k\), and \(f_\chi\) picks
a witness whenever one exists.

Let \(T=\mathrm{Th}(M^{\mathrm{sk}})\).  By
\cref{thm:p03-w03-p03-030-mss}, \(M^{\mathrm{sk}}\) expands to a reduced model
\((M^{\mathrm{sk}},q)\) of a self-extending theory \(T'\) containing (B).
Let \((\mathfrak B,r)\) be its canonical extension.  By
\cref{thm:p03-w03-p03-030-canonical-extension},
\((M^{\mathrm{sk}},q)\preccurlyeq(\mathfrak B,r)\).  The extension is proper
by \cref{thm:p03-w03-p03-030-large-sets}, and it is an end-extension by
\cref{thm:p03-w03-p03-030-end-extension}(1).  Let \(N\) be the \(L\)-reduct of
\(\mathfrak B\).  Then \(M\preccurlyeq N\), \(N\ne M\), and \(N\) end-extends
\(M\).  \(M\) is infinite, having no last element, so \(\abs
N\le\abs M+\aleph_0=\abs M\) by part (5) of the theorem, and \(\abs N\ge\abs
M\).
\end{proof}
\end{annexitem}

\begin{annexitem}{annex:p03-w03-p03-030-pa}{thm:p03-w03-p03-030-pa}{The MacDowell--Specker theorem for Peano arithmetic}
\begin{proof}
All facts below are derived from the seven axioms of \(\mathsf Q\) and
instances of the induction scheme
\(\bigl(\phi(0,\bar p)\wedge\forall x\,(\phi(x,\bar p)\rightarrow
\phi(Sx,\bar p))\bigr)\rightarrow\forall x\,\phi(x,\bar p)\) for
\(\Lang_A\)-formulas \(\phi\) with parameters \(\bar p\), so they hold in
\(\mathcal M\).  Write \(x<y\) for
\(\exists z\,(x+Sz=y)\) and \(x\le y\) for \(x<y\vee x=y\).

\emph{Addition.}  (a) \(0+x=x\), by induction on \(x\): \(0+0=0\), and
\(0+Sx=S(0+x)=Sx\).  (b) \(Sx+y=S(x+y)\), by induction on \(y\):
\(Sx+0=Sx=S(x+0)\), and \(Sx+Sy=S(Sx+y)=S(S(x+y))=S(x+Sy)\).  (c)
\((x+y)+z=x+(y+z)\), by induction on \(z\): both sides are \(x+y\) at
\(z=0\), and \((x+y)+Sz=S((x+y)+z)=S(x+(y+z))=x+S(y+z)=x+(y+Sz)\).  (d)
\(x+y=y+x\), by induction on \(y\): at \(0\) use \(x+0=x\) and (a); and
\(x+Sy=S(x+y)=S(y+x)=Sy+x\) by (b).  (e) \(x+z=y+z\rightarrow x=y\), by
induction on \(z\), using injectivity of \(S\) in the step.

\emph{Order.}  (O1) \(\neg\,x<0\), since \(x+Sz=S(x+z)\ne0\).  (O2)
\(x<Sx\), since \(x+S0=S(x+0)=Sx\).  (O3) \(\neg\,x<x\): from \(x+Sz=x=x+0\),
(d) and (e) give \(Sz=0\), which is false.  (O4) \(<\) is transitive: if
\(x+Sz=y\) and \(y+Sw=u\), then by (c) and the recursion for \(+\),
\(u=x+(Sz+Sw)=x+S(Sz+w)\).  (O5) If \(x<y\), then \(Sx=y\) or \(Sx<y\): let
\(x+Sz=y\).  If \(z=0\), then \(y=Sx\) by (O2)'s computation.  Otherwise
\(z=Sz'\) by the predecessor axiom of \(\mathsf Q\), and
\(y=x+SSz'=S(x+Sz')=Sx+Sz'\) by (b).  (O6) Trichotomy
\(\forall y\,(x<y\vee x=y\vee y<x)\), by induction on \(x\).  For \(x=0\):
either \(y=0\), or \(y=Sy'\) and \(0+Sy'=Sy'\) by (a), so \(0<y\).  For the
step, take \(y\).  If \(y<x\) or \(y=x\), then \(y<Sx\) by (O2) and (O4).  If
\(x<y\), then \(Sx=y\) or \(Sx<y\) by (O5).  (O7) If \(y<Sx\), then \(y\le
x\): otherwise \(x<y\) by (O6), so \(Sx\le y\) by (O5), and with \(y<Sx\) and
(O4) this gives \(y<y\), against (O3).

By (O3), (O4) and (O6), \(<\) is a linear order.  By (O1) and the case
\(x=0\) of (O6), \(0\) is its first element, and by (O2) there is no last
element.

\emph{Both schemes.}  Let \(\mathcal M^<\) be the expansion of \(\mathcal M\)
to \(\Lang_A\cup\{<\}\) by the defined order.  Replacing every atomic
formula \(s<t\) by \(\exists z\,(s+Sz=t)\), with \(z\) fresh, turns each
\(\Lang_A\cup\{<\}\)-formula into an \(\Lang_A\)-formula equivalent to it in
\(\mathcal M^<\).  So induction may be used for every formula
\(\phi(x,\bar p)\) of \(\Lang_A\cup\{<\}\), with parameters.

Least-number principle.  Suppose \(\phi(a)\).  Prove by induction on \(b\)
the formula \(\exists x\,(x\le b\wedge\phi(x))\rightarrow\exists
x\,(\phi(x)\wedge\forall y\,(y<x\rightarrow\neg\phi(y)))\).  At \(b=0\), a
witness \(x\le0\) equals \(0\) by (O1), and \(0\) is least by (O1).  At
\(Sb\), if some \(x\le b\) satisfies \(\phi\), use the hypothesis for \(b\).
Otherwise the witness \(x\le Sb\) is not \(\le b\), so \(x=Sb\) by (O7).
Every \(y<Sb\) satisfies \(y\le b\) by (O7), hence \(\neg\phi(y)\), so
\(Sb\) is least.  The case \(b\defeq a\) gives a least element.

Induction scheme of \cref{def:p03-w03-p03-030-induction-scheme}.  Suppose
\(\phi(a)\) and that every element satisfying \(\phi\) has a larger one
satisfying \(\phi\).  Prove \(\exists y\,(\phi(y)\wedge x<y)\) by induction
on \(x\).  At \(x=0\), take \(y>a\) with \(\phi(y)\); since \(0\le a\), we get
\(0<y\).  At \(Sx\), take \(y>x\) with \(\phi(y)\).  By (O5), either
\(Sx<y\), and \(y\) serves, or \(Sx=y\), and then a larger element
satisfying \(\phi\) serves.

\emph{The extension.}  \(\mathcal M^<\) is a structure for the finite
vocabulary \(\Lang_A\cup\{<\}\) that satisfies the induction scheme and the
least-number principle.  By \cref{thm:p03-w03-p03-030-end-extension-pa} it
has a proper elementary end-extension \(\mathcal N^<\) with
\(\abs{\dom\mathcal N^<}=\abs{\dom\mathcal M}\).  Let \(\mathcal N\) be its
\(\Lang_A\)-reduct.  Then \(\mathcal M\preccurlyeq\mathcal N\) and
\(\mathcal N\ne\mathcal M\).  The sentence \(\forall x\,\forall y\,(x<y
\leftrightarrow\exists z\,(x+Sz=y))\) holds in \(\mathcal M^<\), hence in
\(\mathcal N^<\), so the order of \(\mathcal N^<\) is the order defined in
\(\mathcal N\).  The end-extension property of \(\mathcal N^<\) is therefore
the property stated for \(\mathcal N\).
\end{proof}
\end{annexitem}

\begin{annexitem}{annex:p03-w03-p03-030-like-transfer}{thm:p03-w03-p03-030-like-transfer}{From \(\omega\)-like models to \(\lambda\)-like models}
\begin{proof}
Let \(\Gamma\) be \(\Sigma\) together with the sentence that \(<\) is a linear
order with a first and no last element, the induction scheme and the
least-number principle.  Let \(\Gamma_0\subseteq\Gamma\) be finite, and take a
model of \(\Gamma_0\cap\Sigma\) whose order has type \(\omega\).  In it every
nonempty set has a least element.  A set with an upper bound is finite, so a
nonempty set without a greatest element is unbounded.  So the model
satisfies \(\Gamma_0\).  By the compactness theorem
(\cref{thm:fol-compactness}) \(\Gamma\) has a model, and by the downward
L\"owenheim--Skolem theorem (\cref{thm:downward-ls}) it has a countable
model \(M_0\).

Build \((M_\nu)_{\nu<\lambda}\).  Let \(M_{\nu+1}\) be a proper elementary
end-extension of \(M_\nu\) with \(\abs{M_{\nu+1}}=\abs{M_\nu}\), by
\cref{thm:p03-w03-p03-030-end-extension-pa}.  This applies because
\(M_\nu\) is an elementary extension of \(M_0\), so it satisfies the
induction scheme and the least-number principle, which are sets of
sentences.  At limit stages take unions; by \cref{thm:elementary-chain} each
stage is an elementary substructure of the union.

As in the proof of \cref{thm:p03-w03-p03-030-strongly-like}, \(M_\mu\) is an
initial segment of \(M_\nu\) for \(\mu\le\nu\).  By induction,
\(\abs{M_\nu}\le\abs\nu+\aleph_0\), which is less than \(\lambda\) for
\(\nu<\lambda\) because \(\lambda\) is uncountable.  Let
\(M=\bigcup_{\nu<\lambda}M_\nu\).  Then \(M_0\preccurlyeq M\), so
\(M\models\Sigma\).  Choosing \(a_\nu\in M_{\nu+1}\setminus M_\nu\) gives
\(\lambda\) distinct elements, so \(\abs M\ge\lambda\).  Also
\(\abs M\le\lambda\cdot\lambda=\lambda\).  For \(c\in M\), pick \(\nu\) with
\(c\in M_\nu\).  Then \(\{x:x<c\}\subseteq M_\nu\), since \(M_\nu\) is an
initial segment of \(M\), so this set has size less than \(\lambda\).  Hence
\(M\) is \(\lambda\)-like.
\end{proof}
\end{annexitem}

\begin{annexitem}{annex:p03-w03-p03-030-transfer-mu}{thm:p03-w03-p03-030-transfer-mu}{Transfer from \(\aleph_0\)-like to \(\lambda\)-like models}
\begin{proof}
Let \(\tau\) consist of \(<\) and the symbols occurring in \(\Sigma\), so
\(\abs\tau\le\mu\).  Every sentence of \(\Sigma\) is a \(\tau\)-sentence, so
models may be taken to be \(\tau\)-structures.

\emph{Step 1: the approximations on \(\mathbb N\).}  For each finite
\(s\subseteq\Sigma\) choose a model of \(s\) whose order has type
\(\omega\).  Transport it along the order isomorphism onto
\((\mathbb N,<)\).  This gives a \(\tau\)-structure \(\mathfrak N_s\) with
domain \(\mathbb N\), the usual order, and \(\mathfrak N_s\entails s\).  Fix
a nonprincipal ultrafilter \(\mathcal U\) on \(\mathbb N\); it exists by
\cref{thm:p01-ultrafilter-lemma} applied to the filter of cofinite sets.  It
contains every cofinite set, since it contains no finite set.

\emph{Step 2: least-witness functions.}  Put \(L_0=\tau\), and let
\(L_{k+1}\) add to \(L_k\) a new function symbol \(f_\chi\) for each
\(\mathscr L(Q)\)-formula \(\chi(\bar z,y)\) over \(L_k\), its arity being
the length of \(\bar z\).  Let \(L^*=\bigcup_kL_k\).  There are at most
\(\abs{L_k}+\aleph_0\) formulas over \(L_k\), so \(\abs{L^*}\le\mu\).  For
each \(s\), expand \(\mathfrak N_s\) by recursion on \(k\).  Once the
symbols of \(L_k\) are interpreted, interpret \(f_\chi\) for \(\chi\) over
\(L_k\) as the least \(y\in\mathbb N\) with
\((\mathfrak N_s,\mathcal U)\entails\chi(\bar z,y)\), or as \(0\) if there is
none.  Here the family \(\mathcal U\) interprets \(Q\), and \(\chi\) is
evaluated in the \(L_k\)-structure already built.  Call the resulting
\(L^*\)-structure \(\mathfrak N^*_s\).  A formula over \(L_k\) has the same
truth value in \((\mathfrak N^*_s,\mathcal U)\) as in its \(L_k\)-reduct
with \(\mathcal U\), since satisfaction only uses the symbols occurring in
the formula.

\emph{Step 3: the theory.}  Let \(\Gamma\) be the set of the following
\(\mathscr L(Q)\)-sentences over \(L^*\): the sentence that \(<\) is a linear
order; the universal closures of the four ultrafilter schemes of
\cref{def:p03-w03-p03-030-self-extending-theory} for all formulas over
\(L^*\); the sentences \(\forall\bar z\,(\exists y\,\chi(\bar z,y)\rightarrow
\chi(\bar z,f_\chi(\bar z)))\) for all \(\chi\) over \(L^*\); the universal
closures of all instances of the bounded scheme (B) of
\cref{def:p03-w03-p03-030-scheme-reading} over \(L^*\); and \(\Sigma\).

Let \(\Gamma_0\subseteq\Gamma\) be finite, and let \(s=\Gamma_0\cap\Sigma\).
Then \((\mathfrak N^*_s,\mathcal U)\entails\Gamma_0\).  The order is the
usual one.  The ultrafilter schemes hold because \(\mathcal U\) is an
ultrafilter on \(\mathbb N\) containing no singleton, and every definable
set is a subset of \(\mathbb N\).  Each formula \(\chi\) over \(L^*\) lies
over some \(L_k\), and \(f_\chi\) picks a witness whenever one exists, by
the remark at the end of Step 2.  For (B), fix parameters and
\(z\in\mathbb N\), and suppose that \(\{x:\exists y\,(\phi(x,y)\wedge
y<z)\}\in\mathcal U\).  This set is the union of the finitely many sets
\(\{x:\phi(x,y)\}\) with \(y<z\).  An ultrafilter that contains a finite
union contains one of its members, so \(\{x:\phi(x,y)\}\in\mathcal U\) for
some \(y\), that is, \(\exists y\,Qx\,\phi(x,y)\) holds.  Finally
\(\mathfrak N_s\entails s\), and \(s\) consists of \(\tau\)-sentences.

By \cref{thm:p03-w03-p03-030-compactness}, \(\Gamma\) has a model.  Its
complete \(\mathscr L(Q)\)-theory \(T\) is consistent and contains the
ultrafilter schemes, and it is a Skolem theory, the Skolem term for
\(\chi\) being \(f_\chi(\bar z)\).  So \(T\) is self-extending, and it
contains (B) and \(\Sigma\).

\emph{Step 4: the chain.}  By \cref{thm:p03-w03-p03-030-hull}, \(T\) has a
reduced model \((\mathfrak A_0,q_0)\) with \(\abs{A_0}\le\mu\).  Let
\((\mathfrak A_{\nu+1},q_{\nu+1})\) be the canonical extension of
\((\mathfrak A_\nu,q_\nu)\), and take the unions of
\cref{thm:p03-w03-p03-030-chains} at limit stages \(\nu\le\lambda\).  By that
lemma, each stage is a reduced model of \(T\) and an elementary extension of
every earlier stage.  By \cref{thm:p03-w03-p03-030-end-extension}(1), each
canonical extension is an end-extension, and by
\cref{thm:p03-w03-p03-030-large-sets} it is proper.

By induction on \(\nu\), \(A_\nu\) is an initial segment of \(A_{\nu'}\) for
\(\nu\le\nu'\le\lambda\).  The successor step composes two initial-segment
inclusions.  At a limit \(\nu'\), if \(b\in A_{\nu'}\) and \(b<a\in A_\nu\),
then \(b\in A_\rho\) for some \(\rho\) with \(\nu\le\rho<\nu'\), and \(A_\nu\)
is an initial segment of \(A_\rho\), so \(b\in A_\nu\).

Also by induction, \(\abs{A_\nu}\le\mu+\abs\nu\) for \(\nu<\lambda\).  This
holds at \(0\).  At successors it holds by part (5) of
\cref{thm:p03-w03-p03-030-canonical-extension}, since
\(\abs{L^*}+\aleph_0\le\mu\).  At limits, \(A_\nu\) is a union of \(\abs\nu\)
sets of size at most \(\mu+\abs\nu\).  Since \(\mu<\lambda\) and
\(\abs\nu<\lambda\), every \(A_\nu\) with \(\nu<\lambda\) has fewer than
\(\lambda\) elements.

\emph{Step 5: the model.}  Let \(\mathfrak A\) be the \(\tau\)-reduct of
\(\mathfrak A_\lambda\).  It satisfies \(\Sigma\subseteq T\).  It has at least
\(\lambda\) elements, because each step \(\nu<\lambda\) adds an element.  It
has at most \(\lambda\) elements, being the union of \(\lambda\) sets each of
size less than \(\lambda\).  If \(c\in A_\lambda\), then \(c\in A_\nu\) for
some \(\nu<\lambda\), and \(\{x:x<c\}\subseteq A_\nu\) because \(A_\nu\) is
an initial segment of \(A_\lambda\).  So that set has fewer than
\(\lambda\) elements, and \(\mathfrak A\) is a \(\lambda\)-like model of
\(\Sigma\).
\end{proof}
\end{annexitem}
